% Section 16.4: Naturality | Part: 3 | Depends on: sec:fields, sec:tangent | Status: done | Original lines: 2523-2537
%!TEX root = ../../main.tex
\subsection{Naturality}\label{sec:naturality}

For a smooth map $\phi:M\to N$, a curve $\gamma$ in $M$ and a $1$-form $\omega$ on $N$:
\begin{equation}\label{eq:natural}
\int_{\phi\circ\gamma}\omega=\int_\gamma\phi\pb\omega .
\end{equation}
\emph{Proof in steps:}
\begin{align*}
\int_{\phi\circ\gamma}\omega&=\int_a^b\Bigl\langle\omega_{\phi(\gamma(t))},\frac{\dd}{\dd t}(\phi\circ\gamma)(t)\Bigr\rangle\dd t &&\text{(definition \eqref{eq:lineint})}\\
&=\int_a^b\bigl\langle\omega_{\phi(\gamma(t))},\dd\phi_{\gamma(t)}\dot\gamma(t)\bigr\rangle\dd t &&\text{(chain rule)}\\
&=\int_a^b\bigl\langle(\phi\pb\omega)_{\gamma(t)},\dot\gamma(t)\bigr\rangle\dd t &&\text{(definition \eqref{eq:pullback})}\\
&=\int_\gamma\phi\pb\omega .
\end{align*}
In particular, line integrals of $1$-forms are invariant under changes of coordinates (take $\phi$ a diffeomorphism). In contrast, the classical ``work integral'' $\int\mathbf F\cdot\dd\mathbf r$ is invariant only because it is secretly $\int F^\flat$, where $F^\flat=g(F,\cdot)$ is the covector obtained from $F$ with the Euclidean metric: $F\cdot\dot r=g(F,\dot r)=\langle F^\flat,\dot r\rangle$.
